\chapter{Programas elaborados}
\label{ap:codigo}

La simulación completa (entrenamiento de los modelos, evaluación de los cuatro métodos en Fase~1 y Fase~2, y generación de figuras y tablas) está implementada en Python sobre NumPy, pandas y matplotlib, sin frameworks de aprendizaje automático (PyTorch, TensorFlow, scikit-learn). La red profunda y el Q-Learning se implementan con operaciones matriciales de NumPy; así se evita depender de versiones concretas de esos frameworks.

Se incluyen los extractos del código que implementan cada controlador y la dinámica de la planta.

\lstset{
  language=Python,
  basicstyle=\ttfamily\scriptsize,
  numbers=left, numberstyle=\tiny,
  breaklines=true, columns=fullflexible,
  keepspaces=true, upquote=true,
  frame=single, tabsize=4,
  keywordstyle=\color{UniPrimaryDeep}\bfseries,
  commentstyle=\color{UniMuted}\itshape,
  stringstyle=\color{UniPrimary!80!black},
  showstringspaces=false
}

\begin{lstlisting}[language=Python, caption={Modelo servo reducido del Grupo~1 y controlador PID con anti-windup.}]
# Modelo reducido de primer orden (ZOH, Ts = tau/12) - ecuación 2.3
A = 0.92    # polo discreto = e^(-Ts/tau)
B = 0.08    # ganancia de control = K_f(1-a), K_f = Gp(0) = 1.0
Bd = 0.024  # ganancia de perturbacion de carga = K_d(1-a)
SV_REF = 0.6              # referencia de operacion normalizada
Uff = SV_REF * (1.0 - A) / B  # feedforward: u_ss = 0.60

# y(k+1) = A*y(k) + B*u(k) - Bd*d(k)

# PID discreto con anti-windup condicional
def pid_step(e, st, kp, ki, kd, u_min=0.0, u_max=1.0, aw_lim=6.0, ff=0.0):
    de  = e - st.prev_e
    raw = ff + kp*e + ki*st.integral + kd*de
    u   = float(np.clip(raw, u_min, u_max))
    if not (raw > u_max and e > 0) and not (raw < u_min and e < 0):
        st.integral = float(np.clip(st.integral + e, -aw_lim, aw_lim))
    st.prev_e = e
    return u
\end{lstlisting}

\begin{lstlisting}[language=Python, caption={Controlador difuso Mamdani: fuzzificación, inferencia AND-mínimo y defuzzificación por singletons.}]
# FAM de 35 reglas: 7 conjuntos de error (niveles -3..3) x 5 de derivada (-2..2)
# Consecuente por niveles: c = u_base + nivel_e*0.12 + nivel_de*0.05
NIVEL_E  = {"NL":-3,"NM":-2,"NS":-1,"ZE":0,"PS":1,"PM":2,"PL":3}
NIVEL_DE = {"NB":-2,"NS":-1,"ZE":0,"PS":1,"PB":2}

def velocidad_difusa(e, de, d, u_base):
    mu_e  = membresia_error(e)      # 7 conjuntos
    mu_de = membresia_derivada(de)  # 5 conjuntos
    num = den = 0.0
    for el, me in mu_e.items():
        if me < 1e-9: continue
        for dl, md in mu_de.items():
            w = min(me, md)         # AND = mínimo (Zadeh)
            if w < 1e-9: continue
            c = np.clip(u_base + NIVEL_E[el]*0.12 + NIVEL_DE[dl]*0.05, 0.05, 0.98)
            num += w * c;  den += w
    # Regla adicional de perturbacion: IF d=CON THEN velocidad alta
    dw = membresia_perturbacion(d)["CON"]
    if dw > 1e-9:
        num += dw * np.clip(u_base + 0.30, 0.05, 0.98); den += dw
    return float(np.clip(num / den, 0.05, 0.98)) if den > 1e-9 else float(u_base)
\end{lstlisting}

\begin{lstlisting}[language=Python, caption={Actualización TD(0) del Q-Learning y codificación de estado para la coordinación de esteras.}]
# Estado: (nivel_E1_bin, nivel_E2_bin, perturbacion_bin) -> 9x9x2 = 162 estados
# Bins no uniformes, densos cerca de n* = 1.5:
EDGES = [0.0, 0.5, 0.9, 1.2, 1.4, 1.6, 1.8, 2.1, 2.5, 3.1]
# Acciones: 11 combinaciones de (velocidad, reparto_E1, admision)
ACCIONES = np.array([
    [0.30, 0.50, 1.00],  # lento: dejar que las esteras se llenen
    [0.46, 0.50, 1.00],  # estado estacionario equilibrado (~u_ss)
    [0.58, 0.50, 1.00],  # drenado moderado equilibrado
    [0.70, 0.50, 0.95],  # drenado rapido equilibrado
    [0.84, 0.50, 0.85],  # drenado fuerte
    [0.95, 0.50, 0.78],  # drenado de emergencia
    [0.52, 0.30, 1.00],  # drenado suave + prioridad E2
    [0.52, 0.70, 1.00],  # drenado suave + prioridad E1
    [0.72, 0.30, 0.90],  # drenado rapido + prioridad E2
    [0.72, 0.70, 0.90],  # drenado rapido + prioridad E1
    [0.62, 0.50, 1.00],  # drenado fino cerca del objetivo
])

# Actualización TD(0): regla de Bellman
r = -1.5*iae - 30.0*violacion - 8.0*alarma - 0.20*desbalance
Q[s, a] += alpha * (r + gamma * np.max(Q[s_siguiente]) - Q[s, a])
\end{lstlisting}

\begin{lstlisting}[language=Python, caption={Bucle de simulacion Fase~2: controlador difuso mas supervisor de capacidad.}]
for k in range(1, N_muestras):
    dm = demanda(k)
    d  = perturbacion(k)
    ls = perdida_por_perturbacion(k)     # [0.35*d, 0.175*d]
    nm = max(n[k-1])
    e  = n[k-1] - n_objetivo
    de = e - e_anterior

    # 1. Controlador difuso: velocidad por estera
    uk = [velocidad_difusa(e[i], de[i], d, u_base) for i in range(2)]

    # 2. Supervisor de capacidad (barrera de seguridad)
    if nm > umbral_alarma:
        uk = np.clip(uk + 0.14, 0.05, 1.0)
        admision = 0.72
    elif d > 0:
        uk = np.clip(uk + 0.14, 0.05, 1.0)
        admision = 0.86
    elif nm > umbral_riesgo:
        uk = np.clip(uk + 0.07, 0.05, 1.0)
        admision = 0.90

    # 3. Reparto difuso del robot
    r1 = reparto_difuso(n[k-1,0] - n[k-1,1])

    # 4. Dinamica de las dos esteras (balance de masa; sin recorte superior)
    n[k], salida[k] = paso_estera(n[k-1], uk, r1, admision, dm, ls)
\end{lstlisting}


\begin{lstlisting}[language=Matlab,caption={Configuracion del sistema difuso para el controlador inteligente directo},label={lst:fis_control_directo}]
fis = sugfis('Name','fis_control_directo');

% Entrada e: 7 conjuntos, universo documentado en la Tabla 4.2
fis = addInput(fis,[-1.5 1.5],'Name','e');
e_names = {'NL','NM','NS','ZE','PS','PM','PL'};
e_types = {'trapmf','trimf','trimf','trimf','trimf','trimf','trapmf'};
e_params = {[-1.5 -1.5 -0.7 -0.4],[-0.6 -0.35 -0.12],[-0.22 -0.11 0], ...
            [-0.07 0 0.07],[0 0.11 0.22],[0.12 0.35 0.6],[0.4 0.7 1.5 1.5]};
for i = 1:numel(e_names)
    fis = addMF(fis,'e',e_types{i},e_params{i},'Name',e_names{i});
end

% Entrada de: 5 conjuntos, universo documentado en la Tabla 4.2
fis = addInput(fis,[-0.5 0.5],'Name','de');
de_names = {'NB','NS','ZE','PS','PB'};
de_types = {'trapmf','trimf','trimf','trimf','trapmf'};
de_params = {[-0.5 -0.5 -0.22 -0.10],[-0.18 -0.08 0],[-0.05 0 0.05], ...
             [0 0.08 0.18],[0.10 0.22 0.5 0.5]};
for i = 1:numel(de_names)
    fis = addMF(fis,'de',de_types{i},de_params{i},'Name',de_names{i});
end

% Entrada d: SIN/CON
fis = addInput(fis,[0 1],'Name','d');
fis = addMF(fis,'d','trapmf',[0 0 0.35 0.55],'Name','SIN');
fis = addMF(fis,'d','trapmf',[0.35 0.55 1 1],'Name','CON');

% Salida u_star con consecuentes singleton
fis = addOutput(fis,[0.05 0.98],'Name','u_star');
u_base = 0.60;
e_level = [-3 -2 -1 0 1 2 3];
de_level = [-2 -1 0 1 2];
rules = [];
out_idx = 1;
for ie = 1:numel(e_level)
    for ide = 1:numel(de_level)
        c = min(max(u_base + 0.12*e_level(ie) + 0.05*de_level(ide),0.05),0.98);
        fis = addMF(fis,'u_star','constant',c,'Name',sprintf('C%02d',out_idx));
        rules = [rules; ie ide 0 out_idx 1 1];  % 0 = no importa
        out_idx = out_idx + 1;
    end
end
fis = addMF(fis,'u_star','constant',min(u_base + 0.30,0.98),'Name','PERT');
rules = [rules; 0 0 2 out_idx 1 1];            % regla d=CON -> velocidad alta

% Operadores de inferencia
fis.AndMethod = 'min';
fis.OrMethod = 'max';
fis.DefuzzificationMethod = 'wtaver';

fis = addRule(fis,rules);

writeFIS(fis,'fis_control_directo.fis');
fuzzyLogicDesigner(fis)
\end{lstlisting}

La función siguiente implementa únicamente la barrera dura de capacidad usada en el modelo de coordinación del Capítulo~\ref{ch:arquitectura} (Figura~5.16); los refuerzos de velocidad y los regímenes de riesgo/alarma de las ecuaciones del Capítulo~\ref{ch:supervisor} se implementan en \texttt{supervisor\_inteligente} (Listing~\ref{lst:supervisor_inteligente}) dentro del modelo supervisor del Capítulo~\ref{ch:supervisor}.

\begin{lstlisting}[language=Matlab,caption={Funcion del supervisor de capacidad para la coordinacion R2ET},label={lst:supervisor_capacidad}]
function [u1,u2,r1,alpha] = supervisor_capacidad(u1s,u2s,r1s,n1,n2,nmax)

% Senales propuestas por el controlador primario
u1 = u1s;
u2 = u2s;
r1 = r1s;
alpha = 1;

margen = 0.05;

llena1 = n1 >= (nmax - margen);
llena2 = n2 >= (nmax - margen);

% Bloqueo de alimentacion hacia esteras llenas
if llena1 && ~llena2
    r1 = 0.20;
end

if llena2 && ~llena1
    r1 = 0.80;
end

if llena1 && llena2
    alpha = 0.72;
    r1 = 0.5;
end

% Limites fisicos
if r1 < 0.20
    r1 = 0.20;
elseif r1 > 0.80
    r1 = 0.80;
end

if u1 < 0.05
    u1 = 0.05;
elseif u1 > 1
    u1 = 1;
end

if u2 < 0.05
    u2 = 0.05;
elseif u2 > 1
    u2 = 1;
end

end
\end{lstlisting}

\begin{lstlisting}[language=Matlab,caption={Funcion del supervisor inteligente para ajuste de ganancias, caudal y reparto},label={lst:supervisor_inteligente}]
function [Kp,Ki,Kd,u_plus,alpha,r1,modo] = supervisor_inteligente(n1,n2,d,theta,nmax_lim)

nmax = max(n1,n2);

umbral_riesgo = 0.80*nmax_lim;
umbral_alarma = 0.90*nmax_lim;
umbral_lleno  = nmax_lim - 0.05;

% Valores por defecto: modo normal
Kp = 1.0;
Ki = 1.0;
Kd = 0.0;
u_plus = 0.0;
alpha = 1.0;
modo = 0;

% Seleccion del regimen de operacion
if nmax > umbral_alarma
    Kp = 1.55;
    Ki = 0.65;
    Kd = 0.0;
    u_plus = 0.14;
    alpha = 0.72;
    modo = 2;

elseif d > 0
    Kp = 1.55;
    Ki = 0.65;
    Kd = 0.0;
    u_plus = 0.14;
    alpha = 0.86;
    modo = 2;

elseif nmax > umbral_riesgo
    Kp = 1.25;
    Ki = 0.82;
    Kd = 0.0;
    u_plus = 0.07;
    alpha = 0.90;
    modo = 1;

else
    Kp = 1.0;
    Ki = 1.0;
    Kd = 0.0;
    u_plus = 0.0;
    alpha = 1.0;
    modo = 0;
end

% Reparto base del robot
r1 = 0.5;

if abs(theta) > 0.15
    r1 = 0.5 + 0.8*theta;
end

% Barrera dura de capacidad
if n1 >= umbral_lleno && n2 < umbral_lleno
    r1 = 0.20;
end

if n2 >= umbral_lleno && n1 < umbral_lleno
    r1 = 0.80;
end

if n1 >= umbral_lleno && n2 >= umbral_lleno
    alpha = 0.72;
    r1 = 0.5;
end

if r1 < 0.20
    r1 = 0.20;
elseif r1 > 0.80
    r1 = 0.80;
end

end
\end{lstlisting}

\begin{lstlisting}[language=Matlab,caption={Funcion del PID local de ocupacion para las esteras R2ET},label={lst:pid_ocupacion}]
function [u1,u2,e1,e2] = pid_ocupacion(n_ref,n1,n2,Kp,Ki,Kd,u_plus)

persistent I1 I2 e1_ant e2_ant

if isempty(I1)
    I1 = 0;
    I2 = 0;
    e1_ant = 0;
    e2_ant = 0;
end

Ts = 1;

% Error de ocupacion respecto a la referencia
e1 = n1 - n_ref;
e2 = n2 - n_ref;

I1 = I1 + e1*Ts;
I2 = I2 + e2*Ts;

de1 = (e1 - e1_ant)/Ts;
de2 = (e2 - e2_ant)/Ts;

u1_raw = Kp*e1 + Ki*I1 + Kd*de1 + u_plus;
u2_raw = Kp*e2 + Ki*I2 + Kd*de2 + u_plus;

u1 = u1_raw;
u2 = u2_raw;

if u1 < 0.05
    u1 = 0.05;
elseif u1 > 1
    u1 = 1;
end

if u2 < 0.05
    u2 = 0.05;
elseif u2 > 1
    u2 = 1;
end

e1_ant = e1;
e2_ant = e2;

end
\end{lstlisting}
